\documentclass[12pt,a4paper]{article}
\usepackage[T1]{fontenc}
\usepackage{geometry}
\geometry{margin=2.5cm}
\usepackage{listings}
\usepackage{xcolor}
\usepackage{enumitem}
\usepackage{hyperref}
\usepackage{graphicx}
\usepackage{booktabs}

\definecolor{codegray}{rgb}{0.95,0.95,0.95}
\definecolor{codebg}{rgb}{0.97,0.97,0.97}

\lstdefinestyle{cstyle}{
    language=C,
    backgroundcolor=\color{codebg},
    basicstyle=\ttfamily\small,
    keywordstyle=\color{blue},
    commentstyle=\color{gray!70},
    numbers=left,
    numberstyle=\tiny\color{gray},
    frame=single,
    breaklines=true,
    tabsize=4,
    columns=fullflexible
}

\lstdefinestyle{bashstyle}{
    basicstyle=\ttfamily\small,
    backgroundcolor=\color{codegray},
    frame=single,
    breaklines=true
}

\lstdefinestyle{outstyle}{
    basicstyle=\ttfamily\small,
    backgroundcolor=\color{codebg},
    frame=single,
    breaklines=true
}

\lstset{literate={_}{\_}1}

\title{\textbf{Problem Session 4: Memory Class and API in C Programming}\\
\large CPE 334 Operating Systems, 1/2024}
\author{Department of Computer Engineering, Faculty of Engineering\\
King Mongkut's University of Technology Thonburi}
\date{}

\begin{document}

\maketitle
\vspace{-1cm}

\begin{center}
\textit{Note: This lab must be done on a Unix-based OS (WSL/Ubuntu used)}
\end{center}

\tableofcontents
\newpage

\section{Static Storage Class}

\subsection{Step 1--3: Program static1.c}

\begin{lstlisting}[style=cstyle]
#include <stdio.h>
#include <stdlib.h>

void main()
{
    auto int x = 3;
    while (x > 0)
    {
        static int y = 5;
        y++;
        printf("The value of y is %d\n", y);
        printf("The address of y is %p\n\n", &y);
        x--;
    }
}
\end{lstlisting}

\textbf{Compile and run:}
\begin{lstlisting}[style=bashstyle]
gcc -o static1 static1.c
./static1
\end{lstlisting}

\textbf{Output (run 3 times):}
\begin{lstlisting}[style=outstyle]
===== RUN 1 =====
The value of y is 6
The address of y is 0x62a44b73d010
The value of y is 7
The address of y is 0x62a44b73d010
The value of y is 8
The address of y is 0x62a44b73d010

===== RUN 2 =====
The value of y is 6
The address of y is 0x574453f75010
The value of y is 7
The address of y is 0x574453f75010
The value of y is 8
The address of y is 0x574453f75010

===== RUN 3 =====
The value of y is 6
The address of y is 0x608fe00ee010
The value of y is 7
The address of y is 0x608fe00ee010
The value of y is 8
The address of y is 0x608fe00ee010
\end{lstlisting}

\subsection*{Explanation Step 1--3}

Within a single run, the value of $y$ increases to $6, 7, 8$ and the address of $y$ remains constant across all loop iterations. Normally, a local variable such as \verb|auto int x| is created and destroyed each time the block is entered/exited inside the loop. However, $y$ is \textbf{static}, which means:

\begin{itemize}
    \item It is initialized ($= 5$) only once when the program starts, even though it is declared inside the loop.
    \item Its value and address reside in the \textbf{static memory / data segment}, not on the stack.
    \item On the next loop iteration the value of $y$ is preserved (increasing from $5 \to 6 \to 7 \to 8$).
\end{itemize}

\textbf{About the address between runs:} each time the program is executed the address of $y$ changes (e.g. \texttt{0x62a4...}, \texttt{0x5744...}, \texttt{0x608f...}) because the OS enables \textbf{ASLR (Address Space Layout Randomization)} --- every time the program is exec'ed the OS randomizes where the executable is laid out, so static/global data is placed at a different address in each run for security.

\subsection{Step 4: compile with -no-pie}

\begin{lstlisting}[style=bashstyle]
gcc -no-pie -o static1_nopie static1.c
./static1_nopie
\end{lstlisting}

\textbf{Output (run 3 times):}
\begin{lstlisting}[style=outstyle]
===== RUN 1 =====
The value of y is 6
The address of y is 0x404018
The value of y is 7
The address of y is 0x404018
The value of y is 8
The address of y is 0x404018

===== RUN 2 =====   (same) address = 0x404018
===== RUN 3 =====   (same) address = 0x404018
\end{lstlisting}

\subsection*{Explanation Step 4 (is it different?)}

\textbf{Yes, different} --- with \texttt{-no-pie} the address of $y$ becomes \texttt{0x404018}, constant in every run (no matter how many times it is executed). This is because \texttt{-no-pie} tells the compiler to produce a non-position-independent executable (PIE is disabled), so the program is not affected by ASLR for the code/static segment, placing static data at a fixed address. This is opposite to the PIE (default) case where the address changes every run.

\newpage
\section{Extern Storage Class}

\subsection{Step 1: Program extern1.c (with extern in main)}

\begin{lstlisting}[style=cstyle]
#include <stdio.h>
#include <stdlib.h>

int x = 20;

void display()
{
    extern int x;
    printf(" Display value of x: %d\n", x);
    printf(" Adress of x (in display function): %p\n\n", &x);
}

void main()
{
    extern int x;
    printf(" Print value of x: %d\n", x);
    printf(" Adress of x (in main function): %p\n\n", &x);
    display();
}
\end{lstlisting}

\textbf{Compile and run:}
\begin{lstlisting}[style=bashstyle]
gcc -o extern1 extern1.c
./extern1
\end{lstlisting}

\textbf{Output (run 3 times):}
\begin{lstlisting}[style=outstyle]
===== RUN 1 =====
 Print value of x: 20
 Adress of x (in main function): 0x582057aa2010
 Display value of x: 20
 Adress of x (in display function): 0x582057aa2010
===== RUN 2 =====
 Print value of x: 20
 Adress of x (in main function): 0x6451fd93b010
 Display value of x: 20
 Adress of x (in display function): 0x6451fd93b010
===== RUN 3 =====
 Print value of x: 20
 Adress of x (in main function): 0x557625310010
 Display value of x: 20
 Adress of x (in display function): 0x557625310010
\end{lstlisting}

\subsection*{Explanation Step 1}

Both \texttt{main()} and \texttt{display()} print the value $x = 20$ and use the same address within a given run because $x$ is a \textbf{global variable} declared at global scope (outside any function). \verb|extern| is a declaration meaning ``this variable is defined elsewhere (in another scope)'', so both functions refer to the \textbf{same} variable $x$, which resides in the data segment.

\subsection{Step 2: remove extern in main (Program extern2.c)}

\begin{lstlisting}[style=cstyle]
void main()
{
    // no extern int x;
    printf(" Print value of x: %d\n", x);
    printf(" Adress of x (in main function): %p\n\n", &x);
    display();
}
\end{lstlisting}

\textbf{Output:} identical to \texttt{extern1.c} (value 20, same address in both functions).

\subsection*{Explanation Step 3}

After removing \verb|extern| the result is still identical, because $x$ is a global variable visible from every function automatically. Since main has no local variable named $x$, it uses the global $x$ itself. Adding \verb|extern| in the original code only highlights that we refer to a variable defined externally, but in this case it is redundant because the global is already visible --- hence no difference in the output.

\subsection{Step 4: compile with -no-pie}

\begin{lstlisting}[style=bashstyle]
gcc -no-pie -o extern1_nopie extern1.c
./extern1_nopie
\end{lstlisting}

\textbf{Output:} every run the address of $x$ = \texttt{0x404018} (constant).

\subsection*{Explanation Step 4}

Same as Part 1 --- \texttt{-no-pie} disables ASLR for static/global data, so the global $x$ is placed at a fixed address in every run, unlike the PIE case where the address changes each run.

\newpage
\section{Memory Allocation (malloc/realloc)}

\subsection{Step 1: Program mem1.c (without uncommenting part b)}

\begin{lstlisting}[style=cstyle]
#include <stdio.h>
#include <stdlib.h>

void main()
{
    int *a;
    float *b;
    int c[10];

    printf("\nAddress of Pointer\n");
    printf(">>%p\n", &a);
    printf(">>%p\n", &b);

    printf("\nEffective Address\n");
    printf(">>%p\n", a);
    printf(">>%p\n", b);

    a = (int *) malloc(10*sizeof(int));
    printf("\nAfter malloc Pointer a\n");
    printf(">>%p\n", a);
    printf(">>%p\n", &a[0]);
    printf(">>%p\n", &a[9]);

    printf("\nArray c\n");
    printf(">>%p\n", c);
    printf(">>%p\n", &c[0]);
    printf(">>%p\n", &c[9]);

/*      b = (float*) malloc(10*sizeof(float));
    printf("\nAfter malloc Pointer b\n");
    printf(">>%p\n", b);
    printf(">>%p\n", &b[0]);
    printf(">>%p\n", &b[9]);
*/
    a = (int *) realloc(a, 1000*sizeof(int));
    printf("\nAfter realloc Pointer a\n");
    printf(">>%p\n", a);
    printf(">>%p\n", &a[0]);
    printf(">>%p\n", &a[9]);
    printf(">>%p\n", &a[999]);

    free(a);
//      free(b);
}
\end{lstlisting}

\textbf{Compile and run:}
\begin{lstlisting}[style=bashstyle]
gcc -o mem1 mem1.c
./mem1
\end{lstlisting}

\textbf{Output (Step 2 -- copy into report):}
\begin{lstlisting}[style=outstyle]
Address of Pointer
>>0x7ffd90c86520
>>0x7ffd90c86528

Effective Address
>>(nil)
>>(nil)

After malloc Pointer a
>>0x5fad2e71e2b0
>>0x5fad2e71e2b0
>>0x5fad2e71e2d4

Array c
>>0x7ffd90c86530
>>0x7ffd90c86530
>>0x7ffd90c86554

After realloc Pointer a
>>0x5fad2e71e2b0
>>0x5fad2e71e2b0
>>0x5fad2e71e2d4
>>0x5fad2e71f24c
\end{lstlisting}

\subsection*{Explanation Step 3 (meaning of each part)}

\begin{table}[h]
\centering
\small
\begin{tabular}{@{}p{4.3cm}p{3.6cm}p{8cm}@{}}
\toprule
\textbf{Part} & \textbf{Value} & \textbf{Meaning} \\
\midrule
Address of Pointer \verb|&a| & \verb|0x7ffd90c86520| & address of pointer $a$ itself, on the \textbf{stack} (prefix \texttt{0x7ffd...}); 8 bytes apart from \verb|&b| (64-bit pointer) \\
Address of Pointer \verb|&b| & \verb|0x7ffd90c86528| & address of pointer $b$ on the stack, next to $a$ \\
Effective Address \verb|a| & \verb|(nil)| & what pointer $a$ points to -- NULL because not yet allocated \\
Effective Address \verb|b| & \verb|(nil)| & pointer $b$ points nowhere $\to$ NULL \\
After malloc \verb|a| & \verb|0x5fad2e71e2b0| & pointer $a$ now points to the \textbf{heap} (lower prefix than stack) \\
\verb|&a[0]| & \verb|0x5fad2e71e2b0| & $= a$ (first element of the block) \\
\verb|&a[9]| & \verb|0x5fad2e71e2d4| & offset from \verb|&a[0]| = $9 \times 4$ = 36 bytes (\texttt{int} = 4 bytes) \\
Array $c$ \verb|c| & \verb|0x7ffd90c86530| & array $c$ is on the \textbf{stack} (prefix \texttt{0x7ffd...}) because it is a local array \\
\verb|&c[9]| & \verb|0x7ffd90c86554| & offset $9 \times 4$ = 36 bytes \\
After realloc \verb|a| & \verb|0x5fad2e71e2b0| & \textbf{unchanged} from before -- realloc grown in place \\
\verb|&a[999]| & \verb|0x5fad2e71f24c| & offset from $a$ = $999 \times 4$ = 3996 bytes \\
\bottomrule
\end{tabular}
\end{table}

\textbf{Summary:} local pointers are on the stack, but the memory allocated by malloc is on the heap, and a local array is entirely on the stack. The initial value before malloc is NULL.

\subsection{Step 4: uncomment 2 places (Program mem2.c)}

Uncomment \verb|b = (float*) malloc(...)| and \verb|free(b)|:

\begin{lstlisting}[style=bashstyle]
gcc -o mem2 mem2.c
./mem2
\end{lstlisting}

\textbf{Output (Step 5 -- copy into report):}
\begin{lstlisting}[style=outstyle]
Address of Pointer
>>0x7fff529a64b0
>>0x7fff529a64b8

Effective Address
>>(nil)
>>(nil)

After malloc Pointer a
>>0x604b7e40b2b0
>>0x604b7e40b2b0
>>0x604b7e40b2d4

Array c
>>0x7fff529a64c0
>>0x7fff529a64c0
>>0x7fff529a64e4

After malloc Pointer b
>>0x604b7e40b2e0
>>0x604b7e40b2e0
>>0x604b7e40b304

After realloc Pointer a
>>0x604b7e40b310
>>0x604b7e40b310
>>0x604b7e40b334
>>0x604b7e40c2ac
\end{lstlisting}

\subsection*{Explanation Step 6}

\begin{itemize}
    \item \textbf{Part b:} after malloc $b = \verb|0x604b...2e0|$, \verb|&b[0]| $= \verb|0x604b...2e0|$, \verb|&b[9]| $= \verb|0x604b...304|$ --- block $b$ is on the heap right after block $a$, where \verb|&b[9]| $= b + 9\times4$ (float is 4 bytes).
    \item \textbf{Key difference in the address of $a$ after realloc:}
        \begin{itemize}
            \item In \texttt{mem1} (no allocation of $b$): $a$ after realloc \textbf{stays at the same place} \verb|0x...2e2b0|
            \item In \texttt{mem2} (with $b$ allocated): $a$ after realloc \textbf{changes} to \verb|0x...2b310|
        \end{itemize}
    \item \textbf{Reason:} realloc needs to grow block $a$ from 40 bytes to 4000 bytes (1000 ints)
        \begin{itemize}
            \item In \texttt{mem1}: the free space right after block $a$ is large enough $\to$ realloc grows in place $\to$ address does not change
            \item In \texttt{mem2}: block $b$ was allocated right after $a$, so the space adjacent to $a$ is not enough to grow $\to$ realloc must allocate a new larger block (new address \verb|0x...2b310|), copy the data from the old block, and free the old block $\to$ the address changes
        \end{itemize}
\end{itemize}

\textbf{Result of Step 6:} Yes, different --- especially the address of array $a$ changes because allocating $b$ blocks the in-place growth of block $a$.

\newpage
\section*{Notes}
\begin{itemize}
    \item All steps were executed on the same machine (WSL/Ubuntu) as required by the lab.
    \item The addresses shown (e.g. \verb|0x5fad...|, \verb|0x7ffd...|) may differ each run due to ASLR --- what matters is the pattern/relationship (stack, heap, offset sizes), not the exact numbers.
\end{itemize}

\section*{Submission}
\begin{itemize}
    \item Submit the report via LEB2 before the due date.
    \item This problem session is group work of up to 5 students.
\end{itemize}

\end{document}
